% ---
% id: statistics/binomial-box/prop:rule-exact
% kind: prop
% title: "Properties of the exact binomial quantile rule"
% summary: >-
%   The exact quantiles are the extremal counts meeting their budgets and form
%   an admissible endpoint rule, monotone in the allocation, as do counts
%   decided against upper bounds on the tails. The guard may widen the
%   interval by many counts, yet a moved endpoint stays within 1/n of the
%   honest value.
% prominence: standard
% status: implemented
% uses: [def:interval, def:endpoint-rule, lem:guard, statistics/cor:tail-forms, statistics/lem:exact-strict-tails, statistics/def:conservative-evaluation]
% assumes: []
% refs: []
% provenance: newly-derived
% implements: certiqs_key.statistics.rules:SweepExactRule
% lean:
% ---

\begin{proposition}\label{prop:rule-exact}
Assume a \tref{conservative evaluation}
(Definition~\ref{def:conservative-evaluation}). With the notation of
Definition~\ref{def:interval}, for every $\hat c\in\Chatcon$:
\begin{enumerate}[label=(\roman*),leftmargin=2.2em]
\item the two sets in \eqref{eq:kappaminus}--\eqref{eq:kappaplus} are non-empty
subsets of the finite set $\{0,\dots,n\}$, so $\kappa^{\pm}(\hat c)$ exist and
are unique;
\item \emph{(the exact rule is an admissible endpoint rule)} the pair
$R^{\mp}_{\mathrm{ex}}=\kappa^{\mp}(\hat c)$ satisfies
Definition~\ref{def:endpoint-rule}\,(a)--(c): $\kappa^{\pm}(\hat c)\in
\{0,\dots,n\}$ and
\begin{equation}\label{eq:exact-admissible}
  \Pr\big[\Ncount{\hat c}<\kappa^{-}(\hat c)\big]\le\varepsilon^{-}(\hat c),
  \qquad
  \Pr\big[\Ncount{\hat c}>\kappa^{+}(\hat c)\big]\le\varepsilon^{+}(\hat c);
\end{equation}
if $0<p_{\hat c}<1$ then $\kappa^{-}(\hat c)=0$ at $\varepsilon^{-}(\hat c)=0$
and $\kappa^{+}(\hat c)=n$ at $\varepsilon^{+}(\hat c)=0$; and
$\kappa^{-}(\hat c)=\kappa^{+}(\hat c)=0$ at $p_{\hat c}=0$,
$\kappa^{-}(\hat c)=\kappa^{+}(\hat c)=n$ at $p_{\hat c}=1$. Moreover the
rule is monotone in the allocation: $\kappa^{-}(\hat c)$ is non-decreasing in
$\varepsilon^{-}(\hat c)$ and $\kappa^{+}(\hat c)$ is non-increasing in
$\varepsilon^{+}(\hat c)$. Consequently Lemma~\ref{lem:guard} applies to the
exact rule, and its guarded endpoints $k^{\mp}(\hat c)$ satisfy
\eqref{eq:ordering} and \eqref{eq:budget-met};
\item \emph{(extremality and admissible relaxation)} $\kappa^{-}(\hat c)$ is
the largest and $\kappa^{+}(\hat c)$ the smallest integer in $\{0,\dots,n\}$
satisfying the respective inequality of \eqref{eq:exact-admissible}, and
\emph{every} pair $(\tilde k^{-},\tilde k^{+})$ with
$0\le\tilde k^{-}\le\kappa^{-}(\hat c)$ and
$\kappa^{+}(\hat c)\le\tilde k^{+}\le n$ satisfies \eqref{eq:budget-met} in
place of $k^{\mp}(\hat c)$. Enlarging the interval is therefore always
admissible; shrinking it beyond $[\kappa^{-},\kappa^{+}]$ is not;
\item the defining inequalities are the closed forms
\eqref{eq:tails-used} of Corollary~\ref{cor:tail-forms}, so $\kappa^{\pm}(\hat c)$
are computable from $n$, $p_{\hat c}$ and the allocation alone
(Lemma~\ref{lem:exact-strict-tails});
\item \emph{(degenerate honest mass)} if $p_{\hat c}=0$ then
$k^{-}(\hat c)=k^{+}(\hat c)=0$, and if $p_{\hat c}=1$ then
$k^{-}(\hat c)=k^{+}(\hat c)=n$; in both cases $\varepsilon^{-}(\hat c)$ and
$\varepsilon^{+}(\hat c)$ are unspent, both tails in \eqref{eq:budget-met}
being exactly zero;
\item \emph{(exact cost of the guard for this rule)} write $p=p_{\hat c}$.
Relative to the raw quantile interval
$[\kappa^{-}(\hat c),\kappa^{+}(\hat c)]$, the guard \eqref{eq:guard}
enlarges the interval by exactly \eqref{eq:guard-cost} with
$R^{\mp}=\kappa^{\mp}(\hat c)$, and neither enlargement is bounded by one
count: for every $n\ge1$ and every $p\in(0,1)$, an allocation with
$\varepsilon^{-}(\hat c)\ge1-p^{\,n}$ --- admissible whenever $\ecomp$
permits, since $1-p^{\,n}<1$ --- gives $\kappa^{-}(\hat c)=n$ and a lower
enlargement of $n-\lfloor np\rfloor$ counts, that is, at least $1-p$ in
frequency \emph{uniformly in $n$}; symmetrically
$\varepsilon^{+}(\hat c)\ge1-(1-p)^{n}$ gives $\kappa^{+}(\hat c)=0$ and an
upper enlargement of $\lceil np\rceil$ counts. The exact conditions are
\begin{align}
  \kappa^{-}(\hat c)-k^{-}(\hat c)=0
    &\iff \varepsilon^{-}(\hat c)
      <\Pr\big[\Ncount{\hat c}\le\lfloor np\rfloor\big],
    \label{eq:guard-lo0}\\
  \kappa^{-}(\hat c)-k^{-}(\hat c)\le1
    &\iff \varepsilon^{-}(\hat c)
      <\Pr\big[\Ncount{\hat c}\le\lfloor np\rfloor+1\big],
    \label{eq:guard-lo1}\\
  k^{+}(\hat c)-\kappa^{+}(\hat c)=0
    &\iff \varepsilon^{+}(\hat c)
      <\Pr\big[\Ncount{\hat c}\ge\lceil np\rceil\big],
    \label{eq:guard-up0}\\
  k^{+}(\hat c)-\kappa^{+}(\hat c)\le1
    &\iff \varepsilon^{+}(\hat c)
      <\Pr\big[\Ncount{\hat c}\ge\lceil np\rceil-1\big].
    \label{eq:guard-up1}
\end{align}
Measured instead from the honest value, the guard costs at most $1/n$
unconditionally, \eqref{eq:guard-from-p} of Lemma~\ref{lem:guard}\,(iv);
\item \emph{(evaluation against upper bounds)} let $U^{-}_k$ and $U^{+}_k$,
for $k\in\{0,\dots,n\}$, be rationals determined by $n$, $p_{\hat c}$ and
$k$ alone, with
$U^{-}_k\ge\Pr[\Ncount{\hat c}<k]$, $U^{+}_k\ge\Pr[\Ncount{\hat c}>k]$,
$U^{-}_0=0$ and $U^{+}_n=0$, and with every $U^{\mp}_k$ equal to its tail
when $p_{\hat c}\in\{0,1\}$. Put
\begin{equation*}
  \tilde\kappa^{-}(\hat c)\coloneqq\max\big\{k\in\{0,\dots,n\}:\
    U^{-}_{j}\le\varepsilon^{-}(\hat c)\text{ for }1\le j\le k\big\},
  \qquad
  \tilde\kappa^{+}(\hat c)\coloneqq\min\big\{k\in\{0,\dots,n\}:\
    U^{+}_{k}\le\varepsilon^{+}(\hat c)\big\}.
\end{equation*}
Then $\tilde R^{\mp}=\tilde\kappa^{\mp}(\hat c)$ is an admissible endpoint
rule (Definition~\ref{def:endpoint-rule}\,(a)--(c)), monotone in the
allocation, with $\tilde\kappa^{-}(\hat c)\le\kappa^{-}(\hat c)$ and
$\tilde\kappa^{+}(\hat c)\ge\kappa^{+}(\hat c)$, and it is the exact rule
when every $U^{\mp}_k$ equals its tail. Consequently Lemma~\ref{lem:guard}
applies to it.
\end{enumerate}
\end{proposition}

\begin{proof}
\emph{(i).} $k=0$ belongs to the set in \eqref{eq:kappaminus}, since
$\Pr[\Ncount{\hat c}<0]=0\le\varepsilon^{-}(\hat c)$; and $k=n$ belongs to the
set in \eqref{eq:kappaplus}, since
$\Pr[\Ncount{\hat c}>n]=0\le\varepsilon^{+}(\hat c)$. A non-empty finite set of
integers has a unique maximum and a unique minimum.

\emph{(ii).} \emph{Admissibility.} $\kappa^{-}(\hat c)$ is an element of the
set in \eqref{eq:kappaminus} and $\kappa^{+}(\hat c)$ one of the set in
\eqref{eq:kappaplus}, so both lie in $\{0,\dots,n\}$ and satisfy the
inequalities that define those sets, which are \eqref{eq:exact-admissible}.
This is Definition~\ref{def:endpoint-rule}\,(a).

\emph{Monotonicity in the allocation.} Fix every argument but
$\varepsilon^{-}(\hat c)$. Raising $\varepsilon^{-}(\hat c)$ enlarges the set
in \eqref{eq:kappaminus}, so its maximum does not fall; raising
$\varepsilon^{+}(\hat c)$ enlarges the set in \eqref{eq:kappaplus}, so its
minimum does not rise.

\emph{Zero budget, Definition~\ref{def:endpoint-rule}\,(b).} Let
$0<p_{\hat c}<1$ and write $p=p_{\hat c}$. At $\varepsilon^{-}(\hat c)=0$:
for $k\ge1$, $\Pr[\Ncount{\hat c}<k]\ge\Pr[\Ncount{\hat c}=0]=(1-p)^{n}>0$, so
only $k=0$ lies in the set of \eqref{eq:kappaminus} and
$\kappa^{-}(\hat c)=0$. At $\varepsilon^{+}(\hat c)=0$: for $k\le n-1$,
$\Pr[\Ncount{\hat c}>k]\ge\Pr[\Ncount{\hat c}=n]=p^{n}>0$, so only $k=n$ lies
in the set of \eqref{eq:kappaplus} and $\kappa^{+}(\hat c)=n$.

\emph{Degenerate mass, Definition~\ref{def:endpoint-rule}\,(c).} If
$p_{\hat c}=0$ then $\Ncount{\hat c}=0$ with probability one, so
$\Pr[\Ncount{\hat c}<k]$ is $0$ at $k=0$ and $1$ for $k\ge1$; since
$\varepsilon^{-}(\hat c)<1$ the set in \eqref{eq:kappaminus} is $\{0\}$ and
$\kappa^{-}(\hat c)=0$; and $\Pr[\Ncount{\hat c}>k]=0$ for every $k\ge0$, so
the set in \eqref{eq:kappaplus} is all of $\{0,\dots,n\}$ and
$\kappa^{+}(\hat c)=0$. If $p_{\hat c}=1$ then $\Ncount{\hat c}=n$ with
probability one, so $\Pr[\Ncount{\hat c}<k]=0$ for every $k\le n$ and
$\kappa^{-}(\hat c)=n$; and $\Pr[\Ncount{\hat c}>k]$ is $1$ for $k\le n-1$ and
$0$ at $k=n$, so with $\varepsilon^{+}(\hat c)<1$ the set in
\eqref{eq:kappaplus} is $\{n\}$ and $\kappa^{+}(\hat c)=n$.

The consequence is Lemma~\ref{lem:guard} applied to $R=R_{\mathrm{ex}}$.

\emph{(iii).} Extremality is the definition of the maximum and the minimum
together with the monotonicity of the two tails in $k$
(Lemma~\ref{lem:guard}\,(ii)); the relaxation statement is
Lemma~\ref{lem:guard}\,(ii) with $R^{\mp}=\kappa^{\mp}(\hat c)$.

\emph{(iv).} By Corollary~\ref{cor:tail-forms}, with $p=p_{\hat c}$, the
tails in the inequalities that define the sets of \eqref{eq:kappaminus} and
\eqref{eq:kappaplus} are the closed forms \eqref{eq:tails-used}. Under a
conservative evaluation $p_{\hat c}$, a coordinate of ${\honest}^{-}$ or of
${\honest}^{+}$ (Definition~\ref{def:interval}), is rational, the envelope's
vectors being exact rationals (Definition~\ref{def:conservative-evaluation});
and $\varepsilon^{\mp}(\hat c)$ are rational by
Definition~\ref{def:interval}, which also gives $n\ge1$ and
$p_{\hat c}\in[0,1]$. By
Lemma~\ref{lem:exact-strict-tails}\,(ii) each inequality, one for each
$k\in\{0,\dots,n\}$ in each set, is then decided exactly from $n$,
$p_{\hat c}$, $k$ and $\varepsilon^{\mp}(\hat c)$. So both sets, and with
them their maximum $\kappa^{-}(\hat c)$ and minimum $\kappa^{+}(\hat c)$, are
computed exactly from $n$, $p_{\hat c}$ and the allocation alone.

\emph{(v).} By (ii), $\kappa^{-}(\hat c)=\kappa^{+}(\hat c)=0$ at
$p_{\hat c}=0$, and $\lfloor n\cdot0\rfloor=\lceil n\cdot0\rceil=0$, so
\eqref{eq:guard} gives $k^{-}(\hat c)=k^{+}(\hat c)=0$; at $p_{\hat c}=1$,
$\kappa^{\pm}(\hat c)=n$ and $\lfloor n\rfloor=\lceil n\rceil=n$, so
$k^{-}(\hat c)=k^{+}(\hat c)=n$. In the first case
$\Pr[\Ncount{\hat c}<0]=0$ and $\Pr[\Ncount{\hat c}>0]=0$ because
$\Ncount{\hat c}=0$ almost surely; in the second, $\Pr[\Ncount{\hat c}<n]=0$
and $\Pr[\Ncount{\hat c}>n]=0$ because $\Ncount{\hat c}=n$ almost surely. So
both tails in \eqref{eq:budget-met} are exactly $0$.

\emph{(vi).} The identities are Lemma~\ref{lem:guard}\,(iii) with
$R^{\mp}=\kappa^{\mp}(\hat c)$.

\emph{Unboundedness.} If $\varepsilon^{-}(\hat c)\ge1-p^{\,n}$ then
$\Pr[\Ncount{\hat c}<n]=\Pr[\Ncount{\hat c}\le n-1]=1-p^{\,n}
\le\varepsilon^{-}(\hat c)$, so $k=n$ is admissible in \eqref{eq:kappaminus}
and $\kappa^{-}(\hat c)=n$; the lower enlargement is then
$n-\lfloor np\rfloor\ge n-np=n(1-p)$ counts, i.e.\ at least $1-p$ in
frequency. Such an $\varepsilon^{-}(\hat c)$ is admissible whenever $\ecomp$
permits, because $1-p^{\,n}<1$ for $p\in(0,1)$. Symmetrically, if
$\varepsilon^{+}(\hat c)\ge1-(1-p)^{n}$ then
$\Pr[\Ncount{\hat c}>0]=1-(1-p)^{n}\le\varepsilon^{+}(\hat c)$, so $k=0$ is
admissible in \eqref{eq:kappaplus} and $\kappa^{+}(\hat c)=0$, giving an upper
enlargement of $\lceil np\rceil$ counts.

\emph{The four exact conditions.} We use twice that
$\varepsilon^{\pm}(\hat c)\le\ecomp<1$, which
Definition~\ref{def:interval} records for every admissible input.

First, for every integer $K\ge0$,
\begin{equation*}
  \kappa^{-}(\hat c)\le K
  \quad\Longleftrightarrow\quad
  \Pr\big[\Ncount{\hat c}\le K\big]>\varepsilon^{-}(\hat c).
\end{equation*}
If $K+1\le n$: ($\Rightarrow$) $K+1$ is then not admissible in
\eqref{eq:kappaminus}, by maximality of $\kappa^{-}(\hat c)$, which is exactly
$\Pr[\Ncount{\hat c}<K+1]=\Pr[\Ncount{\hat c}\le K]>\varepsilon^{-}(\hat c)$;
($\Leftarrow$) for every integer $j$ with $K+1\le j\le n$, monotonicity of
$k\mapsto\Pr[\Ncount{\hat c}<k]$ (Lemma~\ref{lem:guard}\,(ii)) gives
$\Pr[\Ncount{\hat c}<j]\ge\Pr[\Ncount{\hat c}\le K]>\varepsilon^{-}(\hat c)$,
so no such $j$ is admissible and the maximum is at most $K$. If instead
$K+1>n$, both sides hold unconditionally: $\kappa^{-}(\hat c)\le n\le K$ by
(i), and $\Pr[\Ncount{\hat c}\le K]=1>\varepsilon^{-}(\hat c)$.
By \eqref{eq:guard-cost} the lower enlargement is $0$ iff
$\kappa^{-}(\hat c)\le\lfloor np\rfloor$ and is at most $1$ iff
$\kappa^{-}(\hat c)\le\lfloor np\rfloor+1$, so taking $K=\lfloor np\rfloor$
and $K=\lfloor np\rfloor+1$ gives \eqref{eq:guard-lo0} and
\eqref{eq:guard-lo1}.

Second, and symmetrically, for every integer $K\le n$,
\begin{equation*}
  \kappa^{+}(\hat c)\ge K
  \quad\Longleftrightarrow\quad
  \Pr\big[\Ncount{\hat c}\ge K\big]>\varepsilon^{+}(\hat c),
\end{equation*}
by the same argument applied to the non-increasing map
$k\mapsto\Pr[\Ncount{\hat c}>k]$ and the minimality of $\kappa^{+}(\hat c)$:
for $K-1\ge0$ the index $K-1$ is inadmissible iff
$\Pr[\Ncount{\hat c}>K-1]=\Pr[\Ncount{\hat c}\ge K]>\varepsilon^{+}(\hat c)$,
and for $K\le0$ both sides hold unconditionally. Taking $K=\lceil np\rceil$
and $K=\lceil np\rceil-1$, both of which are at most $n$, gives
\eqref{eq:guard-up0} and \eqref{eq:guard-up1}.

The distance statement is Lemma~\ref{lem:guard}\,(iv).

\emph{(vii).} Both sets are non-empty: the condition defining the first is
vacuous at $k=0$, and $k=n$ belongs to the second because $U^{+}_n=0$. Write
$\tilde k^{\mp}\coloneqq\tilde\kappa^{\mp}(\hat c)$.
\emph{Admissibility.} If $\tilde k^{-}=0$ then
$\Pr[\Ncount{\hat c}<\tilde k^{-}]=0$; otherwise
$\Pr[\Ncount{\hat c}<\tilde k^{-}]\le U^{-}_{\tilde k^{-}}
\le\varepsilon^{-}(\hat c)$, the last step being the defining condition at
$j=\tilde k^{-}$. And
$\Pr[\Ncount{\hat c}>\tilde k^{+}]\le U^{+}_{\tilde k^{+}}
\le\varepsilon^{+}(\hat c)$. Both endpoints lie in $\{0,\dots,n\}$, so this
is Definition~\ref{def:endpoint-rule}\,(a); it also places $\tilde k^{-}$
in the set of \eqref{eq:kappaminus} and $\tilde k^{+}$ in that of
\eqref{eq:kappaplus}, so $\tilde k^{-}\le\kappa^{-}(\hat c)$ and
$\tilde k^{+}\ge\kappa^{+}(\hat c)$ by the extremality of (iii).
\emph{Zero budget.} Let $0<p_{\hat c}<1$. At $\varepsilon^{-}(\hat c)=0$,
(ii) gives $\kappa^{-}(\hat c)=0$, so $0\le\tilde k^{-}\le0$; at
$\varepsilon^{+}(\hat c)=0$ it gives $\kappa^{+}(\hat c)=n$, so
$n\ge\tilde k^{+}\ge n$. This is Definition~\ref{def:endpoint-rule}\,(b).
\emph{Equality with the exact rule.} If every $U^{\mp}_k$ equals its tail,
the second set is that of \eqref{eq:kappaplus}; and since
$k\mapsto\Pr[\Ncount{\hat c}<k]$ is non-decreasing
(Lemma~\ref{lem:guard}\,(ii)), the condition holds at every $j\le k$
exactly when it holds at $j=k$, so the first set is that of
\eqref{eq:kappaminus}. Hence $\tilde k^{\mp}=\kappa^{\mp}(\hat c)$.
\emph{Degenerate mass.} At $p_{\hat c}\in\{0,1\}$ every bound equals its
tail, so the rule is the exact one there, and (ii) gives
Definition~\ref{def:endpoint-rule}\,(c).
\emph{Monotonicity in the allocation.} The bounds do not depend on the
allocation. Raising $\varepsilon^{-}(\hat c)$ enlarges the first set, so its
maximum does not fall; raising $\varepsilon^{+}(\hat c)$ enlarges the second,
so its minimum does not rise. The consequence is Lemma~\ref{lem:guard}
applied to $R=\tilde R$.
\end{proof}
